\section{A conormal--period argument in numerical dimension two}
\label{nu2:sec:method}

This section gives an independent use of the boundary: an alternative
proof of the following restricted canonical statement.

\begin{theorem}[The curve-base alternative]\label{nu2:thm:alternative}
Let $X$ be a smooth connected projective complex fourfold.  If its
canonical divisor $K_X$ is nef and $\kappa(X,K_X)=0$, then
$\nu(K_X)\ne2$.
\end{theorem}

Here $\nu$ denotes the intersection-theoretic numerical dimension of a
nef divisor.  The effective divisor used below comes from the section
supplied by $\kappa=0$.  No condition on the irregularity or the Euler
characteristic is imposed.  The statement concerns the actual
canonical divisor on a smooth variety; it is not a nonvanishing
assertion for $\kappa=-\infty$ or a statement for arbitrary klt pairs.

The proof has three parts.  First, a dimension-preserving effective
MMP puts the pluricanonical support on a reduced boundary whose
adjoint defines a pencil.  Second, root covers turn the obstruction
to lifting through a nilpotent thickening into a period equation on
that curve.  Positivity of the source line kills its moving part;
a separate trace test kills classes supported at exceptional
parameters.  Finally, finite invariant filtrations convert lifting
into a quadratic lower bound for plurisections, contradicting
$\kappa=0$.

The use of cyclic covers, infinitesimal neighborhoods and Hodge duality
continues a method in Kawamata's alternative proof of the numerical-
dimension-one case for minimal threefolds
\cite[Section~4, Theorem~4.1 and Lemmas~4.2--4.3]{Kawamata1992}.
That result supplies the ancestry of this method, rather than the
fourfold assertion of Theorem~\ref{nu2:thm:alternative}.

\subsection{The reduced boundary and a genuine pencil}\label{nu2:subsec:reduction}

Suppose that $\nu(K_X)=2$.  A nonzero pluricanonical section gives
$0\ne D_X\sim_{\mathbb Q}K_X$, with $D_X\ge0$.  On a log resolution
$p:Y\to X$, let $T$ be its reduced strict-transform support, let $E_Y$
be the reduced exceptional divisor, and write $K_Y=p^*K_X+F$, with
$F\ge0$ exceptional.  Then
\[
 K_Y+T+E_Y\sim_{\mathbb Q}D_Y:=p^*D_X+F+T+E_Y,
 \qquad \operatorname{Supp}D_Y=T+E_Y.
\]
For some rational $c>0$,
$p^*D_X\le D_Y\le(1+c)p^*D_X+F+E_Y$.
Monotonicity, birational invariance and exceptional invariance of
$\kappa$ and Nakayama's $\kappa_\sigma$ therefore give dimensions
$0$ and $2$ for this effective log adjoint
\cite[v1, Lemma~2.5 and proof of Lemma~4.1]{LiuXu2025}.
Effective fourfold minimal-model existence and termination with scaling
give a projective $\mathbb Q$-factorial dlt model
\cite[Corollary~D(i), Lemma~2.9(i)]{LazicTsakanikas2022},
\cite[Theorem~4.1(ii),(iii)]{Birkar2012}:
\begin{equation}\label{nu2:eq:model}
 A=K_V+B\sim_{\mathbb Q}D\ge0,\quad
 B=B_{\mathrm{red}}=\operatorname{Supp}D,\quad
 A\text{ nef},\quad \kappa(V,A)=0,\quad\nu(A)=2.
\end{equation}
Choose the scaling boundary to contain a fixed ample rational divisor.
For a positive limiting scaling parameter, the dlt pair and this
ampleness hypothesis allow the first termination clause; limit zero with every parameter positive uses the
already obtained log minimal model and the second clause; an attained
zero means the current adjoint is nef.  The nonextracting MMP preserves
the displayed support and both dimensions.  The underlying $V$ is klt.

The sheaves $\mathcal O_V$ and $\mathcal O_V(-B)$ are
Cohen--Macaulay by the local vanishing theorem for integral
$\mathbb Q$-Cartier divisors on a klt variety
\cite[Theorem~2 and Example~4.1]{KollarLocal2011}.  Their quotient shows
that $B$ is a Cohen--Macaulay threefold.  Adjunction on the whole union,
including conductor gluing on its normalization, gives an slc pair
$(B,\Delta_B)$ with $K_B+\Delta_B=A|_B$.
Threefold slc abundance makes this restriction semiample
\cite[Theorem~0.1]{Fujino2000}.
Write $D=\sum c_iB_i$, $c_i>0$, and choose an ample divisor $H$.  Then
\[
 0=A^3H=\sum_i c_iA^2B_iH,
 \qquad 0<A^2H^2=\sum_i c_iAB_iH^2.
\]
All summands are nonnegative.  The image of a generated multiple on
each $B_i$ consequently has dimension at most one, and one component
has one-dimensional image.  Two general linear forms with common
center disjoint from this image generate the restricted line bundle
everywhere, including components mapped to points.  After fixing one
divisible integer $q$, we obtain
\begin{equation}\label{nu2:eq:pencil}
 G=\sum_i d_iB_i\sim qA,\quad d_i\in\mathbb Z_{>0},\qquad
 L=\mathcal O_V(G),\qquad
 L|_B\simeq f^*\mathcal O_{\mathbb P^1}(1),
\end{equation}
where $f:B\to\mathbb P^1$ is surjective.  Its fibers need not be
connected.  This is a morphism of the boundary; no ambient fibration
is required.

\subsection{Root covers and the two conormal frames}\label{nu2:subsec:roots}

The pencil exists only on $B$.  We now construct local covers on which its conormal line and absolute dualizing line have compatible frames.  The construction must retain the entire reduced divisor, including components over individual parameter values.

Fix $r$ divisible by $q$ and every $d_i$, and put $a=r/q>0$.
These integers will not vary with the infinitesimal order.
Near a compact reduced fiber of $f$, choose a disk $U$ and a frame
$\ell$ of $\mathcal O(1)|_U$.  There is an analytic neighborhood
$W$ of $B_U$, a root $P^r\simeq L|_W$, and a boundary frame $p$
with $p^r=f^*\ell$.  Indeed choose a simultaneous triangulation compatible with the fiber
and boundary \cite[Theorem~1.10]{Coste2007}.  In successive barycentric
subdivisions take the open stars of the compact fiber subcomplex.
Their regular-neighborhood retractions preserve the boundary
subcomplex.  These stars are cofinal among neighborhoods of the
fiber, giving arbitrarily small $W$ for which $W$ and $W\cap B$
retract compatibly onto it.  The retractions are the elementary
regular-neighborhood consequence of the triangulation, not an
additional assertion of the cited triangulation theorem.
The Kummer obstruction in $H^2(W,\mu_r)$ vanishes because $L$ is
trivial on that fiber.  Twisting the root by a $\mu_r$-torsor,
using the corresponding $H^1$ isomorphisms, gives the prescribed
boundary frame.  Properness then permits shrinking $U$.

Put $Q=\omega_V(B)$ and $M=Q\otimes P^{-a}$, with reflexive powers
understood; $M^{[q]}\simeq\mathcal O_W$.  Normalize the full fiber
product of the cyclic algebras
\[
 \bigoplus_{i=0}^{r-1}P^{-i},\qquad
 \bigoplus_{j=0}^{q-1}M^{[-j]}.
\]
The first multiplication uses the section defining $G$, and the second
uses this specified power trivialization.  Retaining every component
gives a finite cover $\pi:Z\to W$ with group
$\mu_r\times\mu_q$ and quotient $W$.  Its tautological section $y$
has reduced Cartier zero divisor $E$; the ramification over $B_i$ is
$r/d_i$, and the second cover is unramified in codimension one.
The two Hurwitz formulas are
\[
 K_Z=\pi^*\left(K_V+\sum_i(1-d_i/r)B_i\right),\qquad
 K_Z+E=\pi^*(K_V+B).
\]
The finite-map discrepancy rule gives klt $Z$ and lc $(Z,E)$
\cite[Proposition~5.20]{KollarMori1998}.  In this setting it follows by exposing
any divisorial valuation downstairs, normalizing its finite pullback,
and using the DVR Hurwitz coefficient $e-1$:
$A_{\mathrm{up}}(F)=eA_{\mathrm{down}}(E_0)$ for a divisor
$F$ above $E_0$, where $A$ denotes log discrepancy for the
corresponding Hurwitz pairs.  For the second displayed formula
these are $(Z,E)$ and $(V,B)$.  For the first they are $Z$ with
zero boundary and $V$ with boundary
$\sum_i(1-d_i/r)B_i$; the latter is klt because every coefficient
of the reduced dlt boundary has been lowered.  Thus the lc and
klt inequalities hold for all divisorial valuations.
The torsion root gives actual
isomorphisms
\begin{equation}\label{nu2:eq:adjunction}
 \omega_Z(E)\simeq(\pi^*P)^a\simeq\mathcal O_Z(aE),
 \qquad\omega_E\simeq\mathcal O_E(aE).
\end{equation}
Klt Cohen--Macaulayness makes $Z$ Gorenstein; Cartier adjunction
makes $E$ Gorenstein.  The normalized adjunction pair has conductor
coefficient one at nodes and coefficient zero at regular
codimension-one points, so $E$ is slc
\cite[Lemma~11.16 and Theorem~11.17]{Kollar2023}.  As a reduced union of lc
centers it is Du Bois \cite[Theorem~1.4]{KollarKovacs2010}.

The reduced families used here are algebraic, also through exceptional
parameters.  One explicit construction performs the two normalized
cyclic constructions over the root gerbe
$\sqrt[r]{L/V}=[L^\times/\mathbb G_m]$, where
$\lambda$ acts by $\lambda^r$.  An object over a map $h:S\to V$
is a line bundle $P_S$ together with an isomorphism
$P_S^r\simeq h^*L$.  Over $B_U$, the trivial line with power
identification given by $f^*\ell$ defines the selected morphism
$\sigma:B_U\to\sqrt[r]{L/V}|_{B_U}$.  Pull back along $\sigma$
the reduced divisor of the already normalized cover.  Thus the
normalization precedes this restriction. The root-gerbe description is standard
\cite[Definition~2.2 and Proposition~2.3]{AndreiniJiangTseng2015};
on the possibly singular $V$ it follows from the same local line-bundle
trivializations. The normalization and singularity comparisons used
here are given by the following local construction.  Locally the substitution $v=w^r$ in the frame parameter
identifies the construction with a fixed normalized cover times that
parameter; normalization commutes with this étale localization.
It does not require normalization to commute with restriction to $B$.
A specified boundary root gives a finite étale section of the root
gerbe: its pullback along any other root is the finite étale torsor
of power-compatible isomorphisms.  Pulling the reduced divisor back
by this section preserves its total-space singularity properties.
The resulting $E_U\to U$ comes from a projective algebraic family
with reduced pure three-dimensional Gorenstein Du Bois total space.

Globally take a finite cover $\rho:T'\to\mathbb P^1$ with a positive
line $R$ satisfying $R^r=\rho^*\mathcal O(1)$; the degree-$r$ power
map with $T'=\mathbb P^1$ and $R=\mathcal O(1)$ suffices.
The same construction produces a reduced algebraic family
$g':E'\to T'^\circ$ on a nonempty open $T'^\circ$.  Remove branch values, vertical-component
images and the finitely many parameters needed for flatness, slc
surface fibers, coherent base change and topological local systems.
Keep the original families through the removed points.  In what
follows, a good disk, annulus or locus lies inside the specified open
on which these flatness, base-change and local-system properties
hold, or its corresponding open in a local chart family.
An isomorphism between specified boundary roots extends uniquely as
a germ near a compact fiber, since its isomorphism sheaf is a finite
$\mu_r$-torsor and the retractions preserve its sections.  These germ
isomorphisms respect every actual ideal power and its connecting
map.  They give the cocycles comparing all local root choices with
the one family on $T'^\circ$.

Set $I=\mathcal O_Z(-E)$ and $\mathcal A_j=I^j/I^{j+1}$ for
every $j\in\mathbb Z$.  The frame $p$ gives the graded conormal
frame $u=y/p$ of $\mathcal A_1$; no ambient lift of $u$ is chosen.
Under \eqref{nu2:eq:adjunction}, $u^{-a}$ is an absolute dualizing
frame $\Omega$ of $\omega_E$.  A change of base root frame has
$p'=\lambda p$, where $\lambda$ is a nowhere-vanishing holomorphic
function of the base coordinate.  Under this change,
\begin{equation}\label{nu2:eq:frames}
 u'=\lambda^{-1}u,\qquad\Omega'=\lambda^a\Omega.
\end{equation}
At branch parameters the factors comparing an extending frame of $R$
with these root frames are full-disk holomorphic units, bounded with
their inverses on a smaller disk.

\subsection{The Hodge test, including classes supported at a point}\label{nu2:subsec:hodge}

The moving period equation alone cannot detect a class supported over one parameter.  The following test supplies both the growth estimate needed on the punctured curve and detection at the omitted points.

\begin{lemma}[Curve Hodge and supported trace test]\label{nu2:lem:hodge-test}
Let $p:Y\to T$ be projective, where $T$ is a smooth complex
algebraic curve and $Y$ is a reduced, pure three-dimensional,
Cohen--Macaulay Du Bois variety.
Components of $Y$ may map to points.  On an open $T^\circ$ with flat
Du Bois surface fibers, coherent base change and topological local
systems, put $\mathbb H=R^1p_*\mathbb Q$.  Then
\begin{equation}\label{nu2:eq:hodge}
 R^1p_*\omega_Y=\omega_T\otimes F^0\mathcal H^\vee,
 \qquad F^1\mathcal H^\vee=0
 \quad\text{on }T^\circ,
\end{equation}
where $\mathcal H=\mathbb H\otimes\mathcal O$.  For a flat local
section $\alpha$ of $\mathbb H\otimes\mathbb C$, let
$\alpha_{\mathcal O}$ be its image in $R^1p_*\mathcal O_Y$.
The functional corresponding to $s$ is
$\alpha\mapsto\langle s/dt,\alpha_{\mathcal O}\rangle$ under
relative Serre duality.
The variation is graded-polarizable and admissible, with quasi-unipotent
monodromy.  A section $s$ extending over a missing parameter obeys
\begin{equation}\label{nu2:eq:one-pole}
 \langle s/dt,\alpha_{\mathcal O}\rangle
 =O\bigl(|t|^{-1}(1+|\log|t||)^N\bigr)
\end{equation}
in flat frames on bounded-argument sectors.  Moreover, a germ of
$R^1p_*\omega_Y$ supported at a point is zero if its absolute trace
annihilates the image of
$H_c^2(Y_U,\mathbb C)\to H_c^2(Y_U,\mathcal O_Y)$ on one
sufficiently small disk $U$.  Neither assertion at the missing point
requires a flat map or a reduced Du Bois central fiber.
\end{lemma}

\begin{proof}
We verify the filtered comparison first, then the growth estimate and
the supported-class assertion.  Use
increasing filtrations on right $\mathcal D_T$-modules.  The symbol
$\mathbb Q_Y^H$ denotes the unshifted constant object in the derived
category of mixed Hodge modules, with rational realization
$\mathbb Q_Y$; it need not be a single perverse Hodge module.
Write $\mathbb D$ for Verdier duality in that category and set
$N=Rp_*\mathbb Q_Y^H$ and $P=\mathbb DN$.  Let $\mathcal M_i$ be
the right module underlying the $i$th Hodge-module cohomology of $P$.
For this calculation we reset the convention used for the earlier
boundary-graph modules: there is no dimension-dependent Tate twist
of the constant object here.  The lowest index computed below comes
from this unshifted constant object and its dual, and the variations
retain their ordinary decreasing Hodge filtration.  For a point $x\in T$, write $i_x:\{x\}\hookrightarrow T$.
The constant/Du Bois comparison and filtered proper duality
\cite{Saito1990,Saito1989,Schnell2014}, with the Du Bois comparison
\cite[preprint, Corollary~0.3, Section~5.3, (5.3.2), and Section~5.4, (5.4.1)]{Saito2000}, give
\[
 \operatorname{gr}_j^F\operatorname{DR}(P)=0\ (j<0),\qquad
 \operatorname{gr}_0^F\operatorname{DR}(P)=Rp_*\omega_Y[3].
\]
Only finitely many $\mathcal M_i$ occur locally, and their good
filtrations have a common lower bound.  The truncation spectral sequence of coherent cohomology sheaves is
\[
 E_2^{u,v}=\mathcal H^u\bigl(\operatorname{gr}_j^F
                  \operatorname{DR}(\mathcal M_v)\bigr)
 \ \Longrightarrow\ \mathcal H^{u+v}
                  \bigl(\operatorname{gr}_j^F\operatorname{DR}(P)\bigr).
\]
At the lowest possible filtration index, the curve de Rham complex
has only its degree-zero term $F_j\mathcal M_v$.  The sole column
$u=0$ therefore survives unchanged.  The abutment is zero for $j<0$,
so these lowest pieces vanish.  Ascending through the negative grades
repeats this argument; at $j=0$ the abutment in degree $v$ is
$R^{v+3}p_*\omega_Y$.  We obtain
\begin{equation}\label{nu2:eq:lowest}
 F_{<0}\mathcal M_i=0,\qquad
 F_0\mathcal M_i=R^{i+3}p_*\omega_Y.
\end{equation}
No splitting of $P$ is used.  On $T^\circ$, duality turns the
variation of $\mathcal M_{-2}$ into $\mathbb H^\vee(1)$; the right
module shift gives \eqref{nu2:eq:hodge}.  Degree-one proper
cohomology has only types $(0,0),(1,0),(0,1)$, proving $F^1=0$.

Put $\mathcal M=\mathcal M_{-2}$ and fix a coordinate $t$ at an
exceptional parameter $x$.  We use the increasing right-module
$V$-filtration, so $t$ lowers its index and $\partial_t$ raises it.
First, for every $\alpha>0$, strict specializability gives
\[
 F_0\operatorname{gr}_\alpha^V\mathcal M
 =\partial_t\bigl(F_{-1}\operatorname{gr}_{\alpha-1}^V
                    \mathcal M\bigr)=0,
\]
where the last equality follows from \eqref{nu2:eq:lowest}.
The exhaustive, locally discrete $V$-filtration then gives
$F_0\mathcal M\subset V_0\mathcal M$: any positive leading
$V$-grade of an $F_0$ section would contradict the displayed vanishing.
The localized image lies in the Deligne lattice whose flat coefficients
have the form $t^{-1+\lambda}(\log t)^jh(t)$, with
$0\le\lambda<1$ and $h$ holomorphic.  This proves
\eqref{nu2:eq:one-pole}.

At grade zero, the required strictness has a different formulation;
no surjectivity of the exceptional map is assumed.  Write
\[
 \Psi=\bigl(\operatorname{gr}_{-1}^V\mathcal M,
                 F_{\bullet-1}\bigr),\qquad
 \Phi=\bigl(\operatorname{gr}_0^V\mathcal M,F_\bullet\bigr).
\]
These underlie the unipotent nearby and vanishing mixed Hodge modules.
Their canonical morphisms are
\[
 \mathrm{can}=\partial_t:\Psi\longrightarrow\Phi,
 \qquad
 \mathrm{var}=t:\Phi\longrightarrow\Psi(-1).
\]
For increasing filtrations our Tate convention is
$F_p(\mathcal H(k))=F_{p-k}\mathcal H$.  Thus
$F_p\Psi(-1)=F_p\operatorname{gr}_{-1}^V\mathcal M$;
the Tate twist belongs to $\mathrm{var}$, not to $\mathrm{can}$.
The mixed-Hodge-module construction and strictness apply without a
pure decomposition \cite{Saito1990}; the precise filtered formulas
are \cite[v1, (9.3), Sections~20--21, and (30.3)--(30.4)]{Schnell2014}.
Strictness of $\mathrm{can}$ says, for each integer $p$,
\[
 F_p\operatorname{gr}_0^V\mathcal M\cap\operatorname{im}(\partial_t)
 =\partial_t\bigl(F_{p-1}\operatorname{gr}_{-1}^V\mathcal M\bigr).
\]
Taking $p=0$ and again using $F_{-1}\mathcal M=0$ proves
\begin{equation}\label{nu2:eq:strict}
 F_0\operatorname{gr}_0^V\mathcal M\cap
 \operatorname{im}\bigl(\partial_t:
 \operatorname{gr}_{-1}^V\mathcal M\to
 \operatorname{gr}_0^V\mathcal M\bigr)
 =\partial_t F_{-1}\operatorname{gr}_{-1}^V\mathcal M=0.
\end{equation}

We now apply this identity to a point-supported germ
$s\in F_0\mathcal M$.  Some power of $t$ kills $s$.  Its image in
$\operatorname{gr}_0^V\mathcal M$ is nonzero unless $s=0$, since
$t$ is injective on $V_{<0}\mathcal M$.
If $s=u\partial_t$ were a de Rham boundary, consider a nonzero leading
$V$-grade $\beta>-1$ of $u$.  The isomorphism
$\partial_t:\operatorname{gr}_\beta^V\mathcal M\to
\operatorname{gr}_{\beta+1}^V\mathcal M$ would give a nonzero
positive grade of $s$, contradicting $s\in V_0\mathcal M$.
Descending through the finitely many jumps between an index containing
$u$ and $-1$ puts $u$ in $V_{-1}\mathcal M$.  The grade-zero class
of $s$ would then belong to the intersection in
\eqref{nu2:eq:strict}, which is impossible for nonzero $s$.

Equivalently, the local inverse-image complex is
\[
 i_x^*\operatorname{DR}(\mathcal M)\simeq
 \bigl[\operatorname{gr}_{-1}^V\mathcal M
       \xrightarrow{\ \partial_t\ }
       \operatorname{gr}_0^V\mathcal M\bigr],
 \qquad\text{in degrees }-1,0,
\]
with the source and target filtrations of $\Psi$ and $\Phi$
\cite[(30.3)]{Schnell2014}.  A supported $F_0$ germ therefore injects
into $H^0(i_x^*\operatorname{DR}\mathcal M)$.
Finally, the perverse-cohomology spectral sequence is
\[
 E_2^{p,q}=H^p(i_x^*\operatorname{DR}\mathcal M_q)
 \ \Longrightarrow\ H^{p+q}(i_x^*\operatorname{DR}P),
 \qquad p\in\{-1,0\}.
\]
All other columns vanish on a curve, so no higher differential enters
or leaves them.  For $\mathcal M=\mathcal M_{-2}$ its edge sequence is
\[
 0\longrightarrow H^0(i_x^*\operatorname{DR}\mathcal M_{-2})
 \longrightarrow H^{-2}(i_x^*\operatorname{DR}P)
 \longrightarrow H^{-1}(i_x^*\operatorname{DR}\mathcal M_{-1})
 \longrightarrow0.
\]
This supplies the desired injection without splitting $P$.

We specify the comparison with trace, since detection of a supported
class requires the actual pairing.  Let $U$ be a sufficiently small
disk.  The Du Bois comparison identifies the natural map
$\iota:Rp_*\mathbb C_Y\to Rp_*\mathcal O_Y$ with the roof
\[
 \operatorname{DR}(N)\ \xleftarrow{\ \simeq\ }
 F_0\operatorname{DR}(N)\longrightarrow
 \operatorname{gr}_0^F\operatorname{DR}(N).
\]
Filtered transpose duality identifies its transpose with
\begin{equation}\label{nu2:eq:trace-comparison}
 \gamma:Rp_*\omega_Y[3]\simeq
 \operatorname{gr}_0^F\operatorname{DR}(P)
 \simeq F_0\operatorname{DR}(P)
 \longrightarrow\operatorname{DR}(P).
\end{equation}
The middle identification uses the vanishing of the negative graded
pieces, so it is an identification in the filtered derived
comparison, not a choice of a splitting.  Under proper duality and
Verdier duality the induced map is
\[
 H^1(Y_U,\omega_Y)\xrightarrow{\ \gamma\ }
 H^{-2}(U,\operatorname{DR}(P))
 \simeq H_c^2(Y_U,\mathbb C)^\vee,
\]
and its pairing is exactly
\begin{equation}\label{nu2:eq:trace-pairing}
 \langle\gamma(s),\alpha\rangle_{\mathrm{Verdier}}
 =\operatorname{Tr}_{\mathrm{coh}}
       (s\smile\iota(\alpha)).
\end{equation}
This is the compatibility of filtered transpose duality with coherent
and topological trace
\cite[Sections~2.6--2.12 and 3.7--3.13]{Saito1989}.
Shrink $U$ within a constructibility neighborhood of its center.
The preceding supported-germ injection, the two-row perverse-stalk
spectral sequence, and \eqref{nu2:eq:trace-comparison} then identify a
nonzero supported $s$ with a nonzero class in the displayed
finite-dimensional constructible cohomology group.  Verdier duality
detects it by some $\alpha\in H_c^2(Y_U,\mathbb C)$.  Thus annihilating
all the stated coherent images forces $s=0$ by
\eqref{nu2:eq:trace-pairing}.  This proof neither asserts that arbitrary
coherent compact-support cohomology is Hausdorff nor that the map
from constant-sheaf cohomology onto it is surjective.
\end{proof}

The annular form of the same pairing will produce the differential
equation.  On a good annulus $S$ in a simply connected good disk, compact-support
Leray gives
\[
 H_c^2(Y_S,\mathbb C)\twoheadrightarrow
 H_c^1(S,R^1p_*\mathbb C),
\]
with kernel $H_c^2(S,R^0p_*\mathbb C)$.  Its coherent image
factors through $H_c^2(S,p_*\mathcal O_Y)=0$: the annulus is a
Stein curve, and this coherent compact-support group vanishes
above degree one.  Compatibility of the Leray pairings with
trace therefore kills the kernel.  The lift of a circle-dual
class tensored with a flat $\alpha$ consequently tests $s$ by
$\oint\langle s/dt,\alpha_{\mathcal O}\rangle\,dt$, up to one
fixed nonzero normalization.  The test is independent of its lift.

\subsection{The period equation and all-order lifting}\label{nu2:subsec:lifting}

Let $g:E\to U$ be any one of the local full covers above.  Push
$I^j/I^{j+k}$ forward by the underlying continuous map, as sheaves
of complex vector spaces.  This convention does not posit a
holomorphic map from a thickening to $U$.  We prove simultaneously
for every $j\in\mathbb Z$, every $k\ge1$, every parameter point
and every permitted root choice that
\begin{equation}\label{nu2:eq:lifting}
 g_*(I^j/I^{j+k+1})\longrightarrow g_*(I^j/I^{j+k})
 \quad\text{is surjective}.
\end{equation}
At each fixed order neighborhoods may shrink.

Assume the earlier orders.  The current connecting map factors
uniquely through the leading layer as
$\delta_k^j:g_*\mathcal A_j\to R^1g_*\mathcal A_{j+k}$:
the kernel of the leading-layer map lifts by the preceding induction
step.  Shorter truncation images have finite filtrations by coherent
$g_*\mathcal A_i$.  Acyclicity on small disks and annuli gives
leading-section lifts on those entire opens.  The overlap difference
of products has double-error term of order $2k\ge k+1$; hence the collection
$\delta_k=(\delta_k^j)_{j\in\mathbb Z}$ is one
$\mathbb C$-linear graded derivation.
This remains true in negative degrees because $I$ is invertible.
Define $e,b\in R^1g_*\mathcal O_E$ on the full disk by
\[
 \delta_k(t)=e u^k,\qquad\delta_k(u)=b u^{k+1},
 \qquad m=a+k>0.
\]
On the good locus the derivation rule, evaluated in each locally
free fiber, gives $\delta_k(h)=h'e u^k$ for holomorphic $h$.
Consequently
\begin{equation}\label{nu2:eq:obstruction}
 \delta_k(hu^{-m})=(h'e-mhb)\Omega.
\end{equation}

Every value of this last obstruction annihilates the constant-sheaf
tests, on good annuli and on disks through exceptional points.
Indeed its exact sequence is
\[
 0\to\omega_E\to\omega_{(k+1)E}
 \to I^{-a-k}/I^{-a}\to0,
 \qquad\omega_{(k+1)E}=I^{-a-k}/I^{1-a}.
\]
The first map is dual to reduction of structure sheaves.  An
obstruction class maps to zero in the middle term; absolute coherent
trace is functorial for this closed embedding
\cite[\S\S3--5]{RamisRuget1970}.  Since
$\mathbb C_E\to\mathcal O_{(k+1)E}\to\mathcal O_E$ lifts the
constant test, the pairing is zero.  Only the reduced algebraic
family is subjected to Hodge theory; the nilpotent thickening enters
through this analytic duality.

Define functional-valued periods
$\beta_e(\alpha)=\langle e\Omega/dt,\alpha_{\mathcal O}\rangle$
and $\beta_b(\alpha)=\langle b\Omega/dt,\alpha_{\mathcal O}\rangle$.
The annular trace of \eqref{nu2:eq:obstruction} is
\[
 \oint(h'\beta_e-mh\beta_b)\,dt=0.
\]
Functions $h$ may have poles at the omitted center because the
required lifts exist on the annulus.  Integration by parts and the
test $h=(t-c)^{-1}$ give the decisive equation
\begin{equation}\label{nu2:eq:period}
 \nabla\beta_e=-m\beta_b\,dt.
\end{equation}
Write $\dot\lambda=d\lambda/dt$ for a base-frame change as in
\eqref{nu2:eq:frames}.  The derivation gives
\[
 e'=\lambda^ke,\quad b'=\lambda^kb-\lambda^{k-1}\dot\lambda e,
 \quad\beta_e'=\lambda^m\beta_e,\quad
 \beta_b'=\lambda^m\beta_b-\lambda^{m-1}\dot\lambda\beta_e.
\]
Coordinate Jacobians cancel between $\delta_k(t)$ and $\Omega/dt$.
The actual root-germ cocycles therefore glue the obstruction to
\begin{equation}\label{nu2:eq:positive-map}
 R^m|_{T'^\circ}\longrightarrow
 F^0\bigl((R^1g'_*\mathbb Q)^\vee\otimes\mathcal O_{T'^\circ}\bigr).
\end{equation}
Both periods lie in $F^0$, so \eqref{nu2:eq:period} says this map
is killed by the Higgs field.

Before current-order vanishing, $b\Omega$ is already a holomorphic
section of $R^1g_*\omega_E$ on each original full disk.  Comparison
with its algebraic family, valid also for arbitrary holomorphic
frames, gives \eqref{nu2:eq:one-pole} for $\beta_b$.
Radial integration of \eqref{nu2:eq:period} gives
\[
 \beta_e=O\bigl((1+|\log|t||)^{N+1}\bigr).
\]
Finite substitutions and the full-disk unit factors preserve this
bound in a frame of $R^m$ extending over each puncture.  At a branch
point, compare on the good punctured overlap, where $t=t(w)$ is
unramified.  Write $\Omega_t$ for the absolute residue form in
the original family and $\Omega_w$ for its transported class
there.  Then $\Omega_w/dw=(dt/dw)\Omega_t/dt$, whereas the
derivation coefficient satisfies $e_w=(dt/dw)^{-1}e_t$.
The factors cancel.  Thus the bound from the original full disk,
after finite substitution and the full-disk unit comparison of
source frames, introduces no extra pole.

We now use the positivity of the line $R^m$ on the compact curve.
More generally, a Higgs-killed map from a positive line to $F^0$ of a
graded-polarizable rational mixed variation with $F^1=0$ and
quasi-unipotent monodromy is zero when its coefficients have the
logarithmic-power bound in frames of that line extending over every
puncture.  To prove this assertion, suppose the map is nonzero and choose the least $w$ for which the image is generically
in $W_w$; holomorphicity makes this containment global.  Projection
to $\operatorname{gr}_w^W$ is nonzero.  Flatness of $W_w$ preserves
the derivative condition, and a flat basis adapted to $W$ preserves
the coefficient bound.  Equip the chosen pure highest Hodge bundle
with its Hodge metric $H$, and denote its Chern connection by $D$.
For a local image section $s$, let $\pi_s$ be orthogonal projection
onto the line spanned by $s$, away from its zeros.  Highest-step
curvature gives
\[
 \partial_t\partial_{\bar t}\log\|s\|_H^2
 =\frac{\|(1-\pi_s)D_{\partial_t}s\|_H^2}{\|s\|_H^2}\ge0
\]
\cite[Theorem~(5.2), Lemmas~(4.9),(4.13)]{Griffiths1970}.
The incoming Higgs term is absent and the outgoing term vanishes;
no flatness of the image line is asserted.  A constant Tate twist
handles negative Hodge indices.  Schmid's estimate after killing
the finite monodromy part bounds flat-vector Hodge norms by powers
of logarithms \cite[Theorem~(6.6$'$)]{Schmid1973}.
Thus, relative to a positive-curvature metric on the source line,
the logarithmic norm ratio $v$ has
$\sqrt{-1}\partial\bar\partial v\ge\omega_R>0$ and upper growth
$O(\log(1+\log(1/|t|)))$.
Subtract a local potential for $\omega_R$ and then
$\epsilon\log(1/|t|)$.  The maximum principle on shrinking annuli,
followed by $\epsilon\downarrow0$, gives a subharmonic extension
preserving this inequality.  Interior zeros contribute nonnegative
point masses.  Integration over the compact curve contradicts
$\deg R^m>0$.

Perfect relative Serre duality now gives $e=0$ on the good locus;
\eqref{nu2:eq:period} and $m>0$ give $b=0$.  There
$g_*\mathcal O_E$ is locally a product of copies of $\mathcal O_U$:
the primitive idempotents of $R^0g_*\mathbb C$ split the proper
reduced fibers into connected pieces, and base change identifies
each piece's functions with $\mathcal O_U$.  A derivation kills
idempotents and, by $e=0$, curve functions.  It therefore kills
every leading section in every degree.
At an exceptional parameter, each value of $\delta_k^{-m}$ is now
point-supported and annihilates the constant tests.  The Hodge test
above makes it zero.  Applying this to $u^{-m}$ first gives $b=0$;
applying it to $f_0u^{-m}$ for an arbitrary
$f_0\in g_*\mathcal O_E$ gives $\delta_k(f_0)=0$.
This includes functions on vertical components, not only pullbacks
from $U$.  The Laurent-power rule kills every degree and completes
the simultaneous induction \eqref{nu2:eq:lifting}.

\subsection{Finite descent and the quadratic count}\label{nu2:subsec:count}

All conormal lifting obstructions have vanished.  We descend only finite collections of layers and count their global sections.  This avoids any assertion about convergence of a formal splitting.

For every signed integer $j$ define
\[
 J_j=\mathcal O_V\left(-\left\lceil jG/r\right\rceil\right),
 \qquad\mathcal B_j=J_j/J_{j+1},\qquad F_j=f_*\mathcal B_j.
\]
On the full cover with deck group $H$,
$(\pi_*I^j)^H=J_j$: the valuation condition is
$(r/d_i)\operatorname{ord}_{B_i}(h)\ge j$, equivalently
$\operatorname{ord}_{B_i}(h)\ge\lceil jd_i/r\rceil$.
Normality checks these divisorial sheaves in codimension one.
Since $0<d_i/r\le1$, consecutive ceiling differences are zero or
one; hence $\mathcal B_j$ is a coherent sheaf on the reduced $B$.
Exactness of finite pushforward and of finite-group invariants,
and averaging of the lifts in \eqref{nu2:eq:lifting}, give finite
filtrations of $f_*(J_l/J_m)$ with full coherent quotients $F_j$,
for every $l<m$.  The total filtered sheaf is only required to be
a sheaf of complex vector spaces.

Cartier periodicity and \eqref{nu2:eq:pencil} give
$J_{j-r}=J_j\otimes\mathcal O_V(G)$ and $F_{j-r}=F_j(1)$.
For a coherent $Q$ on $\mathbb P^1$, a point outside its torsion
support gives
$0\to Q(b-1)\to Q(b)\to\mathbb C_x^{\operatorname{rk}Q}\to0$;
thus $\chi(Q(b))=\chi(Q)+b\operatorname{rk}Q$ for every integer $b$.
Finite-filtered cohomology is additive, vanishes above degree one,
and satisfies $h^0\ge\chi$.  Consequently
\begin{align}\label{nu2:eq:quadratic-count}
 h^0(V,J_{-nr}/J_0)
 &\ge\sum_{j=-nr}^{-1}\chi(F_j)\notag\\
 &=n\sum_{\ell=0}^{r-1}\chi(F_\ell)
   +\frac{n(n+1)}2\sum_{\ell=0}^{r-1}\operatorname{rk}F_\ell.
\end{align}
Here $\mathcal B_0=\mathcal O_B$ and the surjectivity of $f$ gives
$\operatorname{rk}F_0\ge1$.  Passing through
$0\to\mathcal O_V\to\mathcal O_V(nG)\to J_{-nr}/J_0\to0$
loses at most the fixed number $h^1(V,\mathcal O_V)$.
In particular, the coefficient of $n^2$ in this lower bound is exactly
$\tfrac12\sum_{\ell=0}^{r-1}\operatorname{rk}F_\ell>0$; all remaining
terms are linear or constant.  Thus $h^0(V,\mathcal O_V(nG))$ has positive quadratic growth,
contradicting $G\sim qA$ and $\kappa(V,A)=0$ in
\eqref{nu2:eq:model}.  Every count uses only finitely many
infinitesimal orders; no convergent formal neighborhood or one
neighborhood valid at all orders is needed.  This proves
Theorem~\ref{nu2:thm:alternative}.
